\SourceLeaf{166}
\begin{center}
{\large NACHLASS UND BRIEFWECHSEL.\par}
\medskip
[ÜBER DIE ERSTEN GRÜNDE DER GEOMETRIE.]
\end{center}

\bigskip
\begin{center}\textsc{Gauss} an \textsc{Bolyai}. Helmstedt, 16. December 1799.\end{center}

\medskip
.\, .\, .\, .\, .\, .\, . Es thut mir sehr leid, dass ich unsere ehemalige grössere Nähe nicht benutzt habe, um mehr von Deinen Arbeiten über die ersten Gründe der Geometrie zu erfahren; ich würde mir gewiss dadurch manche vergebliche Mühe erspart haben und ruhiger geworden sein, als jemand wie ich es sein kann, so lange bei einem solchen Gegenstande noch so viel zu desideriren ist. Ich selbst bin in meinen Arbeiten darüber weit vorgerückt (wiewohl mir meine andern ganz heterogenen Geschäfte wenig Zeit dazu lassen); allein der Weg, den ich eingeschlagen habe, führt nicht so wohl zu dem Ziele, das man wünscht und welches Du erreicht zu haben versicherst, als vielmehr dahin, die Wahrheit der Geometrie zweifelhaft zu machen. Zwar bin ich auf manches gekommen, was bei den meisten schon für einen Beweis gelten würde, aber was in meinen Augen so gut wie NICHTS beweist, z.\,B. wenn man beweisen könnte, dass ein geradliniges Dreieck möglich sei, dessen Inhalt grösser wäre als eine jede gegebene Fläche, so bin ich im Stande die ganze Geometrie völlig strenge zu beweisen. Die meisten würden nun wohl jenes als ein Axiom gelten lassen; ich nicht; es wäre ja wohl möglich, dass, so entfernt man auch die drei Endpunkte des Dreiecks im Raume von einander annähme, doch der Inhalt immer unter (infra) einer gegebenen Grenze wäre. Dergleichen Sätze habe ich mehrere, aber in keinem finde ich etwas Befriedigendes. Mach’ doch
